% === Fig 6: amplitude-matched peak-velocity difference (gen_fig6_matched.py) ===
\begin{figure}[t]
\centering
\includegraphics[width=\textwidth]{figures/fig6_matched.pdf}
\caption{\textbf{Amplitude-matched left/right peak-velocity difference $D$ and its association with $A$.}
$D = \tfrac12(l_{\mathrm{leftward}} - l_{\mathrm{rightward}})$ for a leftward and a rightward saccade matched on amplitude and position within a
recording (dimensionless, log units); positive $D$ means a relatively faster left eye in leftward than in rightward saccades. Participant $D$ requires
both R1 sessions and at least 5 matched pairs in total; its repeatability requires at least 5 pairs in each session, hence the smaller $n$ below.
\textbf{a}, Participant $D$ (mean of the two R1 session medians) in the reading (TEX; $n$ = 351, mean .020 [.011, .029])
and random-saccade tasks (RAN; $n$ = 114, mean .039 [.028, .050]); histograms are normalized to density; values beyond the axis range are placed in
the edge bins and counted in the legend. Different TEX and RAN sample sizes do not by themselves establish a task difference. \textbf{b}, Participant $D$ against TEX $A$ (R1 two-session mean, centred mean-speed construction):
TEX $D$, Spearman $\rho$ = .446 [.348, .536] ($n$ = 331); RAN $D$, $\rho$ = .483 [.307, .635]
($n$ = 112); $D$ values below $-$.3 are drawn as triangles at the lower edge, and their number is given in the panel. \textbf{c}, Spearman $\rho$ with 95\% participant-bootstrap CI between single-session TEX $A$ and single-session $D$ (same- and cross-session pairings,
$n$ = 239--255) and, below the line, for the pooled associations of panel b; cohorts and averaging differ between the two groups, so their magnitudes are
compared descriptively. TEX $D$ repeatability (R1 S1--S2):
ICC(A,1) = .319 [.159, .472], $n$ = 193. RAN repeatability was not estimable, because no participant had at least 5 pairs in each session, and the RAN session-specific associations
rest on small samples (Supplementary Table S8).}
\label{fig:matched}
\end{figure}
